%=============================== Marco teórico ================================
\todo{Fundamentos teóricos necesarios para comprender la tesis.}

\section{Conceptos preliminares}
\begin{definicion}\label{def:ejemplo}
Una \emph{definición} se escribe con el entorno \texttt{definicion} y queda
numerada por capítulo (esta es la \cref{def:ejemplo}).
\end{definicion}

Las ecuaciones numeradas se referencian con la \cref{eq:modelo}:
\begin{equation}\label{eq:modelo}
  f(\mathbf{x}) = \sum_{i=1}^{n} w_i\, x_i + b .
\end{equation}

\begin{teorema}\label{teo:ejemplo}
Enunciado de un teorema de ejemplo.
\end{teorema}

\begin{proof}
\todo{Demostración (o quitar el entorno si no aplica).}
\end{proof}

\section{...}
\todo{Desarrollar el resto del marco teórico.}
